% Third round: CPython-level work on fastunknot after 0.2, and how it was measured.
\section{A third round: constant factors in Python, and how to measure them}
\label{sec:round3}

After version 0.2 the Rust work was paused and the Python package was
optimized further, with no change to the mathematics, the verdicts or the
asymptotic behaviour. This section records what was gained, what was tried and
rejected, and an episode about measurement that changed how every number below
was obtained. All results are exact and identical to those of 0.2: 35 unit
tests pass, and \texttt{crosscheck.py} compares the ranks by homological degree
of the new default scanner with the generic 0.2 scanner and the 0.1 set algebra
on random braid closures (400 closures with 2 to 5 strands, no mismatch).

\subsection{What changed}
A profile of 0.2 on the stress closure showed that the cobordism algebra was no
longer the bottleneck: of 1.6 s under the profiler, the scanner's own
bookkeeping (method calls per differential entry, tuple-keyed dictionaries,
hashing of nested frozensets, a heap of pivot candidates) took most of it.
\begin{itemize}
\item \textbf{A scanner written for CPython} (\texttt{scan\_fast.py}). Objects
  of a stage are numbered $0,\dots,N-1$ and live in lists; the objects created
  from one old object form a contiguous block, so a delooping label is an
  offset; matchings are interned as integers, so every cache key is a tuple of
  small integers; transfers of an edge are cached per (source matching, target
  matching, value) for both smoothings at once; cancellation is one inlined
  loop. The generic scanner of 0.2 is kept for the ablation configurations.
\item \textbf{A bucket queue for pivots.} On the stress closure 114\,197 unit
  entries are candidates but only 16\,884 are used; the rest die when a
  neighbour is cancelled. Markowitz costs are tiny, so one list per small cost
  plus a heap for the rare large ones replaces the heap.
\item \textbf{Integer planar geometry} (\texttt{planar.py}). Circles of
  $a\cup\bar b$ are found by walking partners alternately; a crossing is glued
  by following the at most four strands through it; a transfer plan is compiled
  only over the crossing and the circles it touches, the other circles being
  renumbered; and since monomials of different components live on disjoint
  circles, multiplying a bit-packed sum by a monomial $c$ is the shift
  \texttt{<{}< c}.
\item \textbf{Outside the scanner.} The modular Alexander filter builds the
  matrix by sparse rows (at most three nonzero entries per row) and evaluates
  each distinct entry once; the boundary profile of a greedy scan order is
  tracked inside the greedy in $O(1)$ per step and a losing start is abandoned
  early; the Jones frontier scan uses the interned matchings and the
  strand-following gluing. Each was checked to be output-preserving against the
  previous implementation on 500 to 600 random diagrams, including diagrams
  with kinks.
\end{itemize}

\subsection{Results}
Same harness on both sides (\texttt{perf.py}: in one process, best of $N$,
timers around the call only), so the columns are comparable with each other
but \emph{not} with the fresh-process medians of the previous section.
\begin{center}\small
\begin{tabular}{@{}lrrr@{}}
\toprule
Task (milliseconds) & 0.2 & after the scanner work & ratio \\
\midrule
Scan, stress closure (36 crossings, rank 2949) & 575.4 & 216.8 & 2.65 \\
Scan, 41-crossing 4-strand closure & 135.3 & 51.1 & 2.65 \\
Scan, 40-crossing 3-strand closure & 71.0 & 26.2 & 2.71 \\
Scan, 31-crossing 6-strand closure & 59.3 & 21.4 & 2.77 \\
Scan, Conway\#Conway (unfactored) & 114.9 & 41.7 & 2.75 \\
Scan, $T(3,11)$ & 22.8 & 8.7 & 2.61 \\
Scan, Conway knot & 7.7 & 3.0 & 2.54 \\
Scan, chain of 256 kinks & 15.9 & 10.9 & 1.47 \\
Recognize the 8-crossing hard unknot & 1.26 & 0.90 & 1.41 \\
Recognize Conway knot, exact steps only & 7.8 & 3.2 & 2.43 \\
\midrule
Total of the 14 cases of \texttt{perf.py} & 1022.2 & 394.6 & 2.59 \\
\bottomrule
\end{tabular}
\end{center}
The right-hand column is a run whose control loop (see below) spread by 4.8\%.
The changes outside the scanner came later and were measured as paired ratios
new/old with 95\% confidence intervals, which do not need a quiet machine:
recognition of $T(3,61)$ $0.60$ $[0.53, 0.67]$ and then $0.69$ $[0.62, 0.72]$,
together about $0.41$; the chain of 256 kinks $0.78$ $[0.73, 0.82]$; the Jones
scan of the stress closure $0.67$ $[0.46, 0.72]$; recognition of the Conway
knot $0.92$ $[0.89, 0.93]$. Four later sequential runs were all rejected by the
control rule and are not used; they are flagged in \texttt{perf\_log.json}.

\paragraph{A cache that is a trade-off.} The last change is not a free win and
is reported as such. Every crossing brings new edge labels, so plans cached per
stage are never reused across stages, although a long braid compiles the same
few pictures at every crossing: at boundary size 6 there are five matchings,
and about 35 transfer plans were compiled per crossing of $T(3,n)$. Transfer
plans are now also cached under a label-independent key, the matchings and the
crossing relabelled by rank. This is sound because the relabelling is monotone
and circles are numbered by their minimal label whichever matching is walked
first (tested on random matchings). The reuse is 80\% of the plans for
$T(3,11)$, 94\% for $T(3,31)$, 99\% for the chain of kinks, 38\% and 19\% for
the random 3- and 4-strand closures, but 1\% for the stress closure and none
for the Conway knot. Three agreeing A/B runs: scan of $T(3,11)$ $0.61$
$[0.59, 0.66]$, chain of kinks $0.78$, 40-crossing 3-strand closure $0.97$;
\emph{scan of the Conway knot $1.06$ $[1.04, 1.07]$ and recognition of the hard
unknot $1.07$, slower}; no claim on the large scans. Timing partial variants
against the old code (A/A ratio $0.999$) splits the cost into ranking the
boundary at each stage ($+1.5\%$), building and looking up the key ($+3.0\%$)
and storing ($+1.3\%$). Two explanations offered first were wrong: interning
the shapes as integers, against the hashing of nested tuples, changed nothing,
and the ratio was the same with the cyclic garbage collector disabled. An
automatic switch was dropped on evidence: the inputs that profit get their
first hit only after 35 to 150 lookups, while the stress closure gets 12 hits
in its first 64 lookups and then none, so an early hit count predicts the wrong
way. The cache is on by default, because on the corpus the gains are
milliseconds and the losses below 0.2 ms, and \texttt{shape\_cache=False}
restores the previous speed exactly ($1.003$ $[0.999, 1.007]$).

Two remarks. First, the closing sentence of the previous section, that on an
unknot diagram no filter can decide the Python accelerations change nothing,
was true of 0.1 against 0.2 and is no longer true: that case is now 1.4 times
faster, because the scan itself is. Second, the gain is a constant factor:
between 2.5 and 2.8 on every scan in the table except the chain of kinks (1.5),
where fixed work per crossing dominates. Nothing here touches the exponential
growth, and the Rust port remains several times faster still.

\subsection{Measured and not adopted}
All of these were decided with deterministic counters (compositions performed
during cancellation, differential entries written while adding crossings), and
the ranks were identical in every variant.
\begin{itemize}
\item \emph{One queued pivot candidate per source object.} Fewer queue
  operations, but 42\,197 instead of 29\,347 compositions on the stress closure:
  the queued candidate goes stale. Rejected.
\item \emph{Other pivot rules.} First-in-first-out inside a cost bucket needs
  64\% more compositions on the stress closure; the cost
  $(\mathrm{out}-1)+(\mathrm{in}-1)$ is worse on all four inputs tried. The
  product cost with last-in-first-out stays.
\item \emph{Another score for scan orders.} Scoring greedy orders by
  $\sum 2^{|\partial|/2}$ selects the same start as the current rule (maximal,
  then total boundary size) on all four inputs.
\item \emph{Scoring all greedy starts instead of twelve.} On 40 random closures
  the order changed in 18: equal work in 8, less in 6, more in 4, with ratios
  from $0.29$ to $6.69$. The total, $0.952$, is carried by three inputs and
  changes sign without one of them. Not adopted.
\item \emph{Cancelling saddle isomorphisms before writing their entries.}
  Between 74 and 83\% of the entries written are incident to a cancelled pair
  when it is cancelled, but only 30 to 53\% of the cancellations are between
  the two smoothings of one old object, the predictable kind, and they account
  for 24 to 36\% of the entries. That bounds the gain near 15\% on the stress
  closure before the cost of the extra code. Not pursued.
\item \emph{A fast path for single-entry transfers.} It would apply to 13\% of
  the edges of the stress closure: too little to measure. Not written.
\end{itemize}
The experiments on scan orders have a positive content. The score is the weak
point, not the number of candidates: on the stress closure the greedy starts
14 to 16 need 23\,710 compositions and have boundary profile $(8, 226)$, while
start 0, with profile $(8, 230)$, needs 895\,182, that is 38 times more, with
2.6 times the peak size; 20 of the 36 starts do not finish in 4 s. A
matchings-only simulation cannot replace the score, because the gap occurs at
equal boundary size, where there are at most 14 matchings: it is all
multiplicity, i.e.\ homology of the partial tangle, which is what the
computation computes. A cheap predictor of that multiplicity would be worth
more than everything in this section.

\subsection{An episode about measurement}
Midway, a benchmark run showed everything 40\% slower, including code that had
not been touched. No process of ours was running; the desktop's own load had
changed. A fixed pure-Python control loop, added to the harness, went from
5.3 ms to 12 ms \emph{inside one run}. Sequential before/after timing was
therefore abandoned for decisions: \texttt{perf.py} now times the control next
to every case and rejects a run whose control spreads by more than 10\% (eight
of the fifteen logged runs are rejected), and \texttt{ab.py} times the old and
the new code alternately in one process.

The first criterion of \texttt{ab.py}, ``the interquartile range of the paired
ratios lies below 1'', was then caught making false claims in both directions.
A change confined to order selection, which is 0.2 to 1.8\% of the scans
concerned and was measured on its own as 32 to 58\% faster with identical
output, was reported as 14\% slower on one input and 4\% faster on another.
The bound came from asking what share of the run time the change could touch
at all. The tool now times the old code twice per round; on mid-sized scans
this A/A control shows identical code differing by 20 to 30\% between adjacent
runs. A difference is claimed only if the order-statistic 95\% confidence
interval of the median ratio excludes 1 \emph{and} the median lies outside the
interquartile range of the A/A ratios. The one claim that had been published
under the lenient criterion (7 to 13\% from the refinements of the geometry on
small scans, no claim on large ones) was re-tested between the two revisions
under the strict criterion and holds: $0.87$ to $0.92$ with intervals inside
$[0.87, 0.95]$, no claim on the large scans.

A smaller finding of the same kind: the new scanner at first overshot a time
budget of 1 s by up to 0.24 s (the generic scanner: 0.01 s), because a single
crossing was added without looking at the clock. With a check every 512
objects the overshoot is at most 0.03 s on the same runs.
